\section{Schedule independent work within a backend}
\label{sec:arch-scheduling}

The ownership boundary in Section~\ref{sec:arch-ownership-boundary} also
determines where parallel execution is possible. Two queries can read the
same immutable target, but they must not overwrite the same alignment
scratch space. A database-backed reader may expose the right sequence
interface while still requiring access from its creating PostgreSQL thread.
The scheduler therefore needs more than a collection of independent
mathematical problems. It needs inputs that remain readable, separate output
locations, and enough memory for the work that overlaps.

BioSeqDB expresses this work through oneAPI Threading Building Blocks
(oneTBB). A \emph{task} is a unit of executable work, such as processing a
range of independent queries. A \emph{worker thread} is an operating-system
thread that executes such tasks. An \emph{arena} is oneTBB's scheduling
domain with a limit on how many threads may participate simultaneously.
Tasks describe available parallelism without assigning a permanent thread
to each operation. This is the distinction emphasized by oneTBB's
task-oriented and data-parallel interface~\autocite{onetbbBenefits}.

\subsection{PostgreSQL processes and oneTBB threads}
\label{sec:arch-process-boundary}

PostgreSQL executes a client session in a backend process. Its original
server process, traditionally called the \emph{postmaster}, starts backends
and remains responsible for accepting connections. It does not execute
each client's alignment tasks.\footnote{PostgreSQL 17,
\href{https://www.postgresql.org/docs/17/tutorial-arch.html}
{Architectural Fundamentals}.}
Inside an executing backend, BioSeqDB has one process-local shared
scheduling interface. It owns a oneTBB arena used by alignment and other supported
work. The name ``global pool'' refers to this process-local scope.

The extension also registers a separate PostgreSQL background supervisor
when it is preloaded. The supervisor maintains shared scheduling
configuration and service state. It can perform background work such as
servicing index-persistence requests. It does not own one central CPU arena
to which every client sends its sequence pairs.
Figure~\ref{fig:arch-process-arena} separates these roles.

\begin{figure}[htbp]
\centering
\begin{tikzpicture}[x=1cm,y=1cm,font=\footnotesize,
  archbox/.style={draw=black!55,rounded corners=2pt,align=center,
    inner sep=5pt,minimum height=.75cm},
  archarr/.style={-{Stealth[length=2mm]},thick,draw=black!65}]
  \node[archbox,fill=black!5,minimum width=4cm] (postmaster)
    at (6.4,7.1) {PostgreSQL postmaster};
  \node[archbox,fill=orange!10,text width=3.3cm] (supervisor)
    at (1.9,5.7) {BioSeqDB supervisor\\separate background process};
  \node[archbox,fill=orange!6,text width=3.3cm] (shared)
    at (1.9,3.65) {shared configuration\\and batch counters};
  \draw[archarr] (postmaster.west) -|
    node[pos=.75,left] {starts} (supervisor.north);
  \draw[archarr] (supervisor.south) --
    node[right] {publishes} (shared.north);

  \draw[draw=blue!65,rounded corners=3pt,thick,fill=blue!2]
    (4.35,.15) rectangle (13,6.25);
  \node[anchor=west,font=\bfseries,text=blue!65!black]
    at (4.6,5.9) {One query backend process};
  \draw[archarr] (postmaster.east) -| (10.4,6.25);
  \node[anchor=west] at (10.6,6.85) {starts};
  \node[archbox,fill=blue!10,text width=7.1cm] (backend)
    at (8.7,4.7) {Backend main thread\\
      PostgreSQL reads, input preparation, result merge};
  \draw[archarr,dashed] (shared.east) -- (4.75,3.65)
    |- (backend.west);
  \node[align=center,text=black!65] at (2,2.35)
    {Each backend adopts policy.\\Tasks stay in its process.};

  \draw[draw=green!45!black,rounded corners=2pt,fill=green!4]
    (4.85,.65) rectangle (12.5,3.1);
  \node[anchor=west,font=\bfseries,text=green!35!black]
    at (5.1,2.75) {Local oneTBB arena};
  \node[archbox,fill=blue!10,text width=1.9cm] at (6.25,1.65)
    {backend thread\\may execute\\tasks};
  \node[archbox,fill=green!12,text width=1.6cm] at (8.7,1.65)
    {oneTBB\\worker A};
  \node[archbox,fill=green!12,text width=1.6cm] at (11.15,1.65)
    {oneTBB\\worker B};
  \draw[archarr,<->] (backend.south) --
    node[right,align=left] {submit work\\and wait} (8.7,3.1);
\end{tikzpicture}
\caption{The supervisor and the query backend are different PostgreSQL
processes. The arena belongs to the query backend. Its worker threads share
that process's address space, and the backend main thread may participate
in arena computation. Dashed communication carries shared policy, not a
cross-process queue of alignment tasks. The three arena participants are
illustrative.}
\label{fig:arch-process-arena}
\end{figure}

Initialization follows the same boundary. During preload, BioSeqDB
registers PostgreSQL hooks and background-worker metadata. It delays arena
initialization until execution in a backend or the background process.
At a scheduling boundary, the backend adopts the current configuration
generation and configures its own arena. A policy restart can advance that
generation and recreate the caller's arena. Other backends adopt it at a
later synchronization point. This is not an immediate stop and restart of
every query in the database.

oneTBB's arena limit includes the application thread when it participates.
An arena limit of four therefore does not promise four extra workers in
addition to the backend main thread. The scheduling interface's range operations enter
the arena through \texttt{task\_arena::execute}, which permits the caller
to execute work. The configured concurrency is also distinct from a live
census of threads that happen to be busy~\autocite{onetbbTaskArena}.

PostgreSQL's parallel-query workers are different resources. A
\texttt{Gather} plan obtains additional PostgreSQL processes, subject to
PostgreSQL's worker-process limits~\autocite{postgresql17_parallel_resources}.
A oneTBB thread inside a backend does not become one of those processes and
does not consume a PostgreSQL parallel-worker slot. Nevertheless, both
forms of concurrency use the same machine's processors and memory.
BioSeqDB clamps its arena width to one when running inside a PostgreSQL
parallel worker unless composition is explicitly enabled. This is a
defensive policy for such callers. It does not make the alignment custom
scan parallel-safe, nor does it allocate a common CPU quota across all
client sessions. Several ordinary backends can still compete for the same
cores.

\subsection{Choose the unit of parallel work}
\label{sec:arch-work-distribution}

The scheduling interface uses oneTBB range operations, task groups, and
explicit task submission. The biological pipeline's stages remain the
operations described in the earlier algorithm chapters. They are not
assigned to a row of permanent threads, with one thread sketching and
another chaining. Instead, the implementation exposes independent work
where its data and algorithm dependencies permit it.

At dispatch level, eligible in-memory sequence pairs can run independently.
Each item writes to its own result collection. The backend combines those
collections after the range has completed. A \emph{join} waits for submitted
work to finish before its results are combined. A second useful shape is several
queries against one shared target. The backend first validates the query
and target identities and prepares the supported target reader. When the
required resident target can be admitted, queries borrow the same immutable
target state and retained query sketches. They still own separate alignment
scratch and results.

Figure~\ref{fig:arch-query-distribution} shows why the number of tasks need
not equal the number of threads. Query costs vary with the number of anchors
and the extension work they induce. oneTBB can subdivide ranges and schedule
ready tasks as participants become available. Its scheduler favors recently
created local tasks and can steal older ready work from another
participant's deque. This balances available work while seeking locality,
without guaranteeing a particular task order~\autocite{onetbbScheduler}.

\begin{figure}[htbp]
\centering
\begin{tikzpicture}[x=1cm,y=1cm,font=\footnotesize,
  task/.style={draw=blue!60,fill=blue!10,minimum height=.55cm,
    inner sep=2pt,align=center},
  timearr/.style={-{Stealth[length=2mm]},draw=black!60,thick}]
  \node[draw=green!50!black,fill=green!7,rounded corners=2pt,
    minimum width=12.1cm,minimum height=.65cm] (target) at (6.6,5.95)
    {one resident target and immutable index reader};
  \foreach \i in {0,...,7} {
    \pgfmathsetmacro{\xx}{1.2+1.55*\i}
    \node[task,minimum width=1.05cm] (q\i) at (\xx,4.75) {$q_{\i}$};
    \draw[-{Stealth[length=1.5mm]},green!45!black]
      (\xx,5.62) -- (q\i.north);
  }
  \node[text=black!65] at (6.6,3.95)
    {ready query work is assigned dynamically};
  \node[anchor=east] at (2.55,2.95) {backend main};
  \node[anchor=east] at (2.55,1.95) {worker A};
  \node[anchor=east] at (2.55,.95) {worker B};
  \foreach \y in {.95,1.95,2.95}
    \draw[draw=black!25] (2.8,\y) -- (12.4,\y);
  \foreach \lo/\hi/\q in {2.8/4.8/0,4.8/7.5/4,7.5/8.7/7} {
    \draw[fill=blue!13,draw=blue!60] (\lo,2.65) rectangle (\hi,3.25);
    \node at ({(\lo+\hi)/2},2.95) {$q_{\q}$};
  }
  \draw[fill=orange!15,draw=orange!70!black] (2.8,1.65) rectangle (7.4,2.25);
  \node at (5.1,1.95) {$q_1$};
  \foreach \lo/\hi/\q in {2.8/4.3/2,4.3/6.0/3,6.0/7.6/5,7.6/10.6/6} {
    \draw[fill=green!13,draw=green!45!black] (\lo,.65) rectangle (\hi,1.25);
    \node at ({(\lo+\hi)/2},.95) {$q_{\q}$};
  }
  \draw[dashed,black!60] (11.15,.45) -- (11.15,3.45);
  \node[anchor=west,align=left] at (11.35,2.0) {join, then\\merge};
  \draw[timearr] (2.8,.15) -- (11.15,.15);
  \node[anchor=west] at (11.35,.15) {time};
\end{tikzpicture}
\caption{Eight independent queries can share three arena participants.
Each query has its own result slot, regardless of which thread executes it.
Colors distinguish participants. Block lengths and assignments illustrate uneven costs, not measurements or
a guaranteed oneTBB schedule. A range task may also contain several queries.
Shared target bytes remain owned until all their consumers finish.}
\label{fig:arch-query-distribution}
\end{figure}

Within one query, sufficiently large groups of independent alignment
regions can also run in parallel. Large-target recomputation can expose
several admitted segments. These choices depend on the selected algorithm
and on whether the relevant sequence views can safely cross threads.
A PostgreSQL reader does not become shareable because its C++ interface
matches a core concept. Such work must retain a supported resident view or
use the backend's serial reading path, as described in
Section~\ref{sec:arch-ownership-boundary}.

The scheduler provides three modes:
\begin{itemize}
  \item \textbf{Single-thread:} run dispatch and coarse pipeline regions on the caller.
  \item \textbf{Dispatch-parallel:} admit eligible pairs from a batch.
  \item \textbf{Task-parallel:} also allow supported query, region, segment
    and index-fragment work.
\end{itemize}
These are runtime scheduling policies,
not another dimension of biological algorithm selection. Running a
supported region serially does not by itself choose a different minimizer,
chainer, or aligner.

Coarse parallel regions share a process-wide guard. Once one such region
has acquired it, another coarse region executes serially instead of opening
another level of parallel work. For example, a parallel query batch does not also
multiply its active workers by independently parallelizing every query's
alignment regions. The guard applies to these coarse regions, rather than
prohibiting all nested oneTBB operations. Workload thresholds also avoid
submitting very small ranges whose scheduling overhead could exceed their
useful work.

Each participating thread needs its own installed computation context for
temporary pipeline storage. For context-aware pipeline ranges, the pool
supplies contexts scoped to the parallel region and reuses an already
installed context when possible.
The results outlive temporary scratch through their own owners.
Consequently, work stealing changes which thread executes ready work, but
does not grant a task access to another thread's mutable alignment scratch
storage.

\subsection{Admit memory before increasing overlap}
\label{sec:arch-admission}

A CPU limit alone cannot bound memory use. Four simultaneous target-index
builds can require roughly four times the temporary storage of one build,
even if each task uses only one processor. BioSeqDB therefore combines the
arena's concurrency limit with a process-local memory-admission controller.
Admission means obtaining a reservation, represented by a move-only permit,
before starting work whose retained allocation is covered by that permit.

The controller uses the common memory policy's view of the process
footprint and available host memory. It also applies an optional index-memory
cap and accounts for bytes already promised to live permits. Reservations
are updated atomically, so simultaneous tasks in the same process cannot
each treat the same remaining reservation capacity as their own. A rejected
request returns without waiting for another task to free memory. The caller
can reduce the admitted wave, use a supported serial or streaming path, or
report that the selected route cannot proceed.

A \emph{wave} is the group of work items whose input and temporary storage
may overlap. It is smaller than the logical batch when the whole batch
would retain too much state. Preparation belongs inside this bound. Waiting
until after every sequence has been copied would leave the copies outside
the intended memory limit. Conversely, a lazy PostgreSQL source need not be
flattened merely because its locator belongs to a batch. The selected
reader and compute path determine which bytes must become resident.

\begin{figure}[htbp]
\centering
\begin{tikzpicture}[x=1cm,y=1cm,font=\footnotesize]
  \node[anchor=east] at (2.35,3.3) {permit};
  \node[anchor=east,align=right] at (2.35,2.35) {wave input\\and scratch};
  \node[anchor=east] at (2.35,1.4) {task A};
  \node[anchor=east] at (2.35,.65) {task B};
  \draw[fill=orange!14,draw=orange!70!black] (2.6,3.0) rectangle (10.7,3.6);
  \node at (6.65,3.3) {capacity reserved};
  \draw[fill=blue!12,draw=blue!65] (4.1,2.0) rectangle (10.15,2.7);
  \node at (7.1,2.35) {storage remains owned};
  \draw[fill=green!13,draw=green!45!black] (5.2,1.1) rectangle (7.7,1.7);
  \draw[fill=green!13,draw=green!45!black] (5.2,.35) rectangle (8.8,.95);
  \draw[dashed,black!60] (9.05,.15) -- (9.05,2.85);
  \node[anchor=south,align=center] at (9.05,3.95) {join and\\merge};
  \draw[-{Stealth[length=1.5mm]},black!60] (9.05,3.91) -- (9.05,3.64);
  \node[anchor=south,align=center] at (2.6,3.95) {reserve};
  \node[anchor=south,align=center] at (4.65,3.95) {prepare};
  \node[anchor=south,align=center] at (11.4,3.95) {release storage,\\then permit};
  \draw[-{Stealth[length=2mm]},black!65,thick] (2.6,-.1) -- (12.9,-.1);
  \node[anchor=west] at (11.5,.55) {next wave};
  \draw[-{Stealth[length=2mm]},black!65] (10.8,.65) -- (11.35,.65);
\end{tikzpicture}
\caption{An admitted wave retains its input and temporary storage through the
join. Completing one task is insufficient to release bytes another task
still reads. In the ordinary permit-retaining path shown here, covered
storage is released before its reservation becomes available to subsequent
work. Output that must survive the wave has a separate owner.}
\label{fig:arch-wave-lifetime}
\end{figure}

Figure~\ref{fig:arch-wave-lifetime} distinguishes reservation from
allocation. A permit records admitted capacity. It does not own the bytes
or prove that the operating system has reserved physical pages. The
ordinary release order is therefore to destroy the covered allocation and
then return its permit. Paths that reconcile an allocation with the measured
process footprint can adjust reservation accounting, but still need an
owner for the bytes and their readers.

Algorithm~\ref{alg:arch-wave-scheduling} summarizes the CPU scheduling
boundary. It deliberately leaves biological implementation selection
outside the loop. An unsupported requested implementation is handled by
its admission policy, rather than silently replaced because fewer CPU
tasks can run.

\begin{algorithm}[htbp]
\caption{CPU work under backend ownership and bounded admission}
\label{alg:arch-wave-scheduling}
\small
\begin{algorithmic}[1]
  \Require supported CPU request and a bounded batch
  \State Backend adopts its scheduling configuration
  \For{each next row or query wave}
    \State Validate source and index identities on the backend
    \State Obtain admission for required resident input and scratch
    \If{mode, work size, ownership, and admission allow parallel work}
      \State Prepare retained, worker-shareable input
      \State Run independent items in the local arena into separate result slots
      \State Join the submitted range
    \Else
      \State Execute the supported serial or streaming path
    \EndIf
    \State Check PostgreSQL interrupts at backend-controlled boundaries
    \State Merge completed results into their continuing owners
    \State Release the wave's input and temporary storage
    \State Return remaining admission
  \EndFor
\end{algorithmic}
\end{algorithm}

These decisions bound controlled overlap within one process. They do not
provide a cluster-wide memory reservation. Another backend can allocate
after a host-memory sample, and unavailable host probes cannot establish
how much physical memory remains. Persistent results and other database
state also have their own lifetimes. The useful architectural guarantee is
the explicit ownership and admission of controlled work, not a claim that
any number of sessions or any input can never exhaust memory.

\subsection{Completion, cancellation, and progress}
\label{sec:arch-scheduler-completion}

Submission does not transfer PostgreSQL's API to a worker. Relation access,
PostgreSQL allocation, and interrupt checks stay on the backend main thread.
The worker runs the permitted computation and returns results or a failure
to that thread. For explicitly tracked asynchronous tasks, the owner registers
cleanup resources before submission. A bounded error record carries worker
failure, and completion is published before joining can release the borrowed
storage. The CustomScan lifetime rules in
Section~\ref{sec:customscan-batch-state-machine} supply the enclosing
statement cleanup. They need not be implemented again by every numeric
algorithm.

This boundary limits cancellation responsiveness. The ordinary parallel
range completes before the backend resumes its per-row interrupt checks.
The native aligner does not continuously call PostgreSQL's cancellation
API while computing an alignment. The explicitly tracked task wrapper can
observe a cancellation request before invoking its callback, but that is
not evidence of interruption at every point inside an already running
kernel. A join protects ownership even when the caller no longer needs a
result. It does not imply that already executing CPU work stops instantly.

The scheduler must also retain enough runnable work to make progress.
oneTBB is intended primarily for computation that does not spend long
periods blocked on IO or mutexes~\autocite{onetbbBlocking}. Nonblocking memory
admission and the coarse-region guard avoid important forms of nested
waiting and resource multiplication. They do not prove general deadlock
freedom. BioSeqDB's input path, for example, can overlap parsing with
backend writes through a bounded queue, so a claim that every pool task is
nonblocking would be incorrect. Its producer and consumer require their
own close-and-join protocol. The common arena supplies execution resources,
while each operation retains responsibility for a valid dependency graph
and bounded ownership.
